% !TEX root = hw4-answers.tex
\chead{\textbf{Exercises week 4}}

\begin{document}
\flushleft\textbf{Topic: Markov kernels}

\section{Composition of Markov Kernels}
\begin{enumerate}
\item $\Zcal, \Tcal, \Wcal$ are assumed to be finite discrete spaces with Markov kernels:\\
 
  \[Q(Z|W):\, \Wcal \dshto \Zcal,\quad (C,w) \mapsto Q(Z \in C|W=w),  \]
  \[K(W|T):\, \Tcal \dshto \Wcal,\quad (D,t) \mapsto K(W \in D|T=t),  \]

and corresponding stochastic matrices:

\[\tilde Q = 
\begin{bmatrix}
0.1 & 0.7 & 0.2 \\
0.3 & 0   & 0.4 \\
0.6 & 0.3 & 0.4 \\
\end{bmatrix} 
\qquad
\tilde K = 
\begin{bmatrix}
0.2 & 0.1 & 0.6 \\
0.3 & 0.1 & 0 \\
0.5 & 0.8 & 0.4 \\
\end{bmatrix}\]

    
Compute the stochastic matrix $\tilde P$ for the composed  Markov kernel: $P(Z|T):=Q(Z|W) \circ K(W|T)$.
(1pt)
\begin{gradedsolution}
$\tilde P = \begin{bmatrix}
0.33 & 0.24 & 0.14 \\
0.26 & 0.35 & 0.34 \\
0.41 & 0.41 & 0.52 \\
\end{bmatrix}$.
\end{gradedsolution}

\item Consider two Markov kernels:
  \[K(W|T):\, \Tcal \dshto \Wcal,\quad (D,t) \mapsto K(W \in D|T=t),  \]
  \[Q(Z|W):\, \Wcal \dshto \Zcal,\quad (C,w) \mapsto Q(Z \in C|W=w),  \]
  that have linear multivariate Gaussian densities (w.r.t.\ Lebesgue measure):
\begin{align*} 
k(w|t) &= \Ncal(w|A \cdot t +b, \Sigma),  \\
q(z|w) & = \Ncal(z| C \cdot w + d, \Gamma).
\end{align*}

Compute the density of the composition: $P(Z|T):=Q(Z|W) \circ K(W|T)$. First, recall how the general formula for the density of a composition is computed.

Suggested steps:
\begin{enumerate}
  \item Write the log joint density $\log p(z, w|t)$ as a quadratic polynomial in $\R^{\dim \Wcal + \dim \Zcal}$ using joint precision matrix $\Lambda$:
  \[ 
    -\frac{1}{2} \mat{z \\ w}^T  \mat{\Lambda_{zz} & \Lambda_{zw} \\ \Lambda_{zw}^T & \Lambda_{ww}}\mat{z \\ w} + \mat{b_1 \\ b_2}^T\mat{z \\ w} + \text{const}
  \]
  \item Invert to find the covariance matrices:
  \[ 
    -\frac{1}{2} \mat{z \\ w}^T  \mat{\Sigma_{zz} & \Sigma_{zw} \\ \Sigma_{zw}^T & \Sigma_{ww}}^{-1}\mat{z \\ w} + \mat{b_1 \\ b_2}^T\mat{z \\ w} + \text{const}
  \]
  % \item Invert to find the joint precision matrix $\Lambda$:
  % \[ 
  %   -\frac{1}{2} \mat{z \\ w}^T \Lambda \mat{z \\ w} + \mat{b_1 \\ b_2}^T\mat{z \\ w} + \text{const}
  % \]
  \item Complete the square to obtain Gaussian form:
  \[ 
    \mat{z-\mu_z \\ w-\mu_w}^T \mat{\Sigma_{zz} & \Sigma_{zw} \\ \Sigma_{zw}^T & \Sigma_{ww}}^{-1}\mat{z-\mu_z \\ w-\mu_w} + \text{const}
  \]
  % \item Invert to bring into Gaussian form:
  % \[ 
  %   % \mat{z-\mu_z \\ w-\mu_w}^T \mat{\Xi_{zz}^{-1} & \Xi_{zw}^{-1} \\ {\Xi_{zw}^{-1}}^T & \Xi_{ww}^{-1}}\mat{z-\mu_z \\ w-\mu_w} + \text{const}
  %   \mat{z-\mu_z \\ w-\mu_w}^T \Xi^{-1}\mat{z-\mu_z \\ w-\mu_w} + \text{const}
  % \]

  \item Read off the marginal $p(z|t)$.
\end{enumerate}
You can use the following identity for symmetric matrices $M,N$ and matrix $T$ without proof:
\begin{align*}
\mat{N &-NT \\ -T^TN  &M+T^TNT}^{-1} = \mat{N^{-1}+TM^{-1}T^T&  TM^{-1}  \\M^{-1}T^T & M^{-1} }
\end{align*}
% You can use the following general identities without proof, for any matrices $A,B,C,D$:
% \begin{enumerate}
%   \item If $B$ symmetric:
%   \[x^T B^{-1} x -2b^Tx = (x-Bb)^TB^{-1}(x-Bb) + \text{const in $x$}\]
%   \item For a 2x2 block matrix, with $M=(A - BD^{-1}C)^{-1}$
%   \[\mat{A & B \\ C & D}=\mat{M & -MBD^{-1} \\ -D^{-1}CM & D^{-1}+D^{-1}CMBD^{-1}}^{-1}\]
%   \item \[(A+CBC^T)^{-1}=A^{-1}-A^{-1}C(B^{-1}+C^TA^{-1}C)^{-1}C^TA^{-1}\]
% \end{enumerate}
% $\textit{Hint: compute}$ $p(z|t) = \int_\Wcal q(z|w) \; k(w|t)dw $.
(1pt)

\end{enumerate}

\begin{gradedsolution}
  Define $\mu=At+b$.
\begin{align*}
&(z-Cw-d)^T \Gamma^{-1}(z-Cw-d)+(w-\mu)^T\Sigma^{-1}(w-\mu) \\
  &=\mat{z \\w}^T \mat{\Gamma^{-1} & -\Gamma^{-1}C \\ -C^T \Gamma^{-1} & \Sigma^{-1} + C^T \Gamma^{-1}C}\mat{z \\ w} -2\mat{z \\ w}^T\mat{\Gamma^{-1}d \\ \Sigma^{-1}\mu - C^T\Gamma^{-1}d} + \text{const}  \\
  &=\mat{z \\w}^T \mat{\Gamma + C \Sigma C^T & C\Sigma \\ \Sigma C^T & \Sigma}^{-1}\mat{z \\ w} -2\mat{z \\ w}^T\mat{\Gamma^{-1}d \\ \Sigma^{-1}\mu - C^T\Gamma^{-1}d} + \text{const}  \\
  % &=\mat{z \\w}^T \mat{\Sigma_{zz} & \Sigma_{zw} \\ \Sigma_{zw}^T & \Sigma_{ww}}^{-1}\mat{z \\ w}^{-1} -2\mat{z \\ w}^T\mat{d \\ \Sigma^{-1}\mu} + \text{const}
\end{align*}
Completing the square is:
\[-\frac{1}{2} x\Sigma^{-1}x+b^Tx=-\frac{1}{2}(x-\Sigma b)^T\Sigma^{-1}(x-\Sigma b)+\text{const}
\]
So the mean is:
\begin{align*}
\mat{\mu_z \\ \mu_w} =\mat{\Gamma + C \Sigma C^T & C\Sigma \\ \Sigma C^T & \Sigma}\mat{\Gamma^{-1}d \\ \Sigma^{-1}\mu - C^T\Gamma^{-1}d}=\mat{C\mu+d \\ \mu}
\end{align*}
and we can read off that:
\[p(z) = \mathcal{N}(z|C \mu + d, \Gamma + A \Sigma A^T)= \mathcal{N}(z|C At + Cb + d, \Gamma + A \Sigma A^T)\]

% \begin{align}
% p(z|t) &= \int_\Wcal q(z|w) \; k(w|t)dw \\
% &= \int_\mathcal{W} \frac{1}{\sqrt{2\pi}} \;  \exp\left(-\frac{1}{2} (w-(At+b))^T \Sigma^{-1}(w-(At+b))\right) \; \\ & \frac{1}{\sqrt{2\pi}}  \exp\left(-\frac{1}{2} (z-(Cw+d))^T \Gamma^{-1}(z-(Cw+d))\right) dw \\
% \end{align} 
\end{gradedsolution}
\section{Markov kernels} %(Bonus Exercise -- you don't need to hand this in)}
\begin{enumerate} 
  \item (1 pt) 
    Let $f: \mathcal{X} \to \mathcal{Y}$ be any function and $\mathcal{B}_{\mathcal{X}}$ a $\sigma$-algebra on $\mathcal{X}$.
    Additionally, let $\mathcal{B}$ be any set of subsets of $\mathcal{Y}$, not necessarily a $\sigma$-algebra.
    Assume that $f$ is $\mathcal{B}_{\mathcal{X}}-\mathcal{B}$--measurable, with which we mean that for each $B \in \mathcal{B}$, we have $f^{-1}(B) \in \mathcal{B}_{\mathcal{X}}$.
    Show that $f$ is then already $\mathcal{B}_\mathcal{X}-\mathcal{B}_{\mathcal{Y}}$--measurable, where $\mathcal{B}_{\mathcal{Y}} \coloneq \sigma{(\mathcal{B})}$ is the smallest $\sigma$-algebra containing $\mathcal{B}$.

    The bottom line, to be used in the second part of the exercise, is the following: measurability can be tested on sets which generate the $\sigma$-algebra of the codomain.
    
    \emph{Hint: Proceed as follows: Show that $\mathcal{B} \subseteq f_*\mathcal{B}_{\mathcal{X}}$. 
    Deduce that $\mathcal{B}_{\mathcal{Y}} \subseteq f_* \mathcal{B}_{\mathcal{X}}$.
  Deduce from that, in turn, that $f$ is $\mathcal{B}_{\mathcal{X}}-\mathcal{B}_{\mathcal{Y}}$-measurable.}
\item (2 x 1 pt)
    Show that a Markov kernel $K$ induces a measurable mapping:
    \[ \bar K:\, \Tcal \to \Pcal(\Wcal),\quad  \bar K(t)(A)  \coloneq K(A|t).  \]% t \mapsto K(X|T=t) = \lp A \mapsto K(X \in A|T=t) \rp\]
    Vice versa, show that any measurable mapping $\mathcal{K}: \mathcal{T} \to \mathcal{P}(\mathcal{W})$ induces a Markov kernel
    \begin{equation}
      \widetilde{\mathcal{K}}: \mathcal{B}_{\mathcal{W}} \times \mathcal{T} \to [0,1],\quad \widetilde{\mathcal{K}}(D \mid t) \coloneq \mathcal{K}(t)(D).
      \label{eq:other_direction_def}
    \end{equation}
    For doing this exercise, recall Definition 4.2.10 
     the $\sigma$-algebra of  $\Pcal(\Wcal)$ (see lecture notes).\\
\end{enumerate}

\begin{gradedsolution}
  \begin{enumerate}
    \item Recall that
      \begin{equation}
	f_* \mathcal{B}_{\mathcal{X}} = \{A \subseteq \mathcal{Y} \mid f^{-1}(A) \in \mathcal{B}_{\mathcal{X}}\}.
	\label{eq:recall_definition}
      \end{equation}
      Since we have $f^{-1}(B) \in \mathcal{B}_{\mathcal{X}}$ for all $B \in \mathcal{B}$ by assumption, we obtain $\mathcal{B} \subseteq f_* \mathcal{B}_{\mathcal{X}}$.
      Now, since $f_* \mathcal{B}_{\mathcal{X}}$ is a $\sigma$-algebra containing $\mathcal{B}$, and since $\mathcal{B}_{\mathcal{Y}}$ is by definition the smallest $\sigma$-algebra containing $\mathcal{B}$, we obtain $\mathcal{B}_{\mathcal{Y}} \subseteq f_* \mathcal{B}_{\mathcal{X}}$.
      By definition of $f_* \mathcal{B}_{\mathcal{X}}$, again, this means that for all $B \in \mathcal{B}_{\mathcal{Y}}$ we have $f^{-1}(B) \in \mathcal{B}_{\mathcal{X}}$.
      This, in turn, means that $f$ is $\mathcal{B}_{\mathcal{X}}-\mathcal{B}_{\mathcal{Y}}$-measurable, proving the claim.
    \item 
      \begin{enumerate}
	\item Assume that $K: \mathcal{B}_{\mathcal{W}} \times \mathcal{T} \to [0,1]$ is a Markov kernel. 
	  Define $\overline{K}$ as in the exercise. 
	  If $t \in \mathcal{T}$, then $\overline{K}(t) \in \mathcal{P}(\mathcal{W})$ according to Definition 6.1.3 i), which shows well-definedness of the map $\overline{K}: \mathcal{T} \to \mathcal{P}_{\mathcal{W}}$.
	  We need to show that it is also measurable, i.e. that preimages of measurable sets are measurable.
	  Note that the $\sigma$-algebra on $\mathcal{P}(\mathcal{W})$ is generated by sets of the form $\text{ev}_D^{-1}(A)$, where $D \subseteq \mathcal{W}$ is measurable, $A \subseteq [0,1]$ is measurable, and $\text{ev}_D: \mathcal{P(\mathcal{W})} \to [0,1]$ is an evaluation map, according to Definition 4.2.10.
	  For these generating sets, we obtain:
	  \begin{align*}
	    \overline{K}^{-1}(\text{ev}_D^{-1}(A)) & = \{ t \in T \mid \overline{K}(t) \in \text{ev}_D^{-1}(A)\} \\
	    & = \{t \in T \mid \text{ev}_D(\overline{K}(t)) \in A\} \\
	    & = \{t \in T \mid \overline{K}(t)(D) \in A\} \\
	    & = \{t \in T \mid K(D \mid t) \in A\} \\
	    & = K(D \mid \bullet)^{-1}(A).
	    \label{eq:derivation_measurable}
	  \end{align*}
	Since according to Definition 6.1.3 ii) the map $K(D \mid \bullet): \mathcal{T} \to [0,1]$ is measurable, this shows the measurability of $\overline{K}^{-1}(\text{ev}_D^{-1}(A))$.
	Using part 1 and that the sets $\text{ev}_D^{-1}(A)$ generate the $\sigma$-algebra on $\mathcal{P}(\mathcal{W})$ proves the measurability of $\overline{K}$.

      \item Note that $\widetilde{\mathcal{K}}(\bullet \mid t) = \mathcal{K}(t) \in \mathcal{P}(\mathcal{W})$, which proves property i) of Definition 6.1.3.
      Now, let $D \in \mathcal{B}_{\mathcal{W}}$ arbitrary and $t \in \mathcal{T}$.
      We obtain
      \begin{align*}
	\widetilde{\mathcal{K}}(D \mid t) & = \mathcal{K}(t)(D) \\
	& = \text{ev}_D(\mathcal{K}(t)) \\
	& = (\text{ev}_D \circ \mathcal{K})(t),
      \end{align*}
      which shows $\widetilde{K}(D \mid \bullet) = \text{ev}_D \circ \mathcal{K}: \mathcal{T} \to [0,1]$.
      Since $\mathcal{K}$ is measurable by assumption and $\text{ev}_D$ is measurable by the definition of the $\sigma$-algebra on $\mathcal{P}(\mathcal{W})$ given in Definition $4.2.10$, and since compositions of measurable functions are measurable, we obtain the measurability of $\widetilde{\mathcal{K}}(D \mid \bullet)$.
      This shows that $\widetilde{\mathcal{K}}$ is a Markov kernel.
	\end{enumerate}
    \end{enumerate}
\end{gradedsolution}



\end{document}
